\documentclass[11pt]{article}
\usepackage[margin=1in]{geometry}
\usepackage{booktabs}
\usepackage{array}
\usepackage{amsmath}
\usepackage{siunitx}
\usepackage{hyperref}

\title{Torsion-Initialized Newton Band Parameter Study}
\author{}
\date{2026-07-12}

\begin{document}
\maketitle

\section{Purpose}

This report studies how the torsion-initialized Newton band parameters
\(\alpha_{T,1}\), \(\alpha_{T,2}\), \(\beta_{\Phi,1}\), and
\(\beta_{\Phi,2}\) affect the agreement between the torsion-designed density
band and the final converged equilibrium density band in
\texttt{scripts/dolfinx\_torsion\_initialized\_newton.py}.

\section{Numerical Setup}

Two fixed smooth-star Gmsh meshes were used.  The initial bracket was run on
\[
\texttt{run\_outputs/strategyA\_band\_study\_20260712/fixed\_mesh/smooth\_star\_h007\_n300.msh},
\]
which has \(16{,}238\) triangles and \(130{,}505\) scalar degrees of freedom
for \(p=4\).  The later parameter sweeps and all saved final PyVista figures
were run on the coarser fixed mesh
\[
\texttt{run\_outputs/strategyA\_band\_study\_20260712/fixed\_mesh/smooth\_star\_h010\_n220.msh},
\]
which has \(9{,}110\) triangles and \(73{,}321\) scalar degrees of freedom for
\(p=4\).  The default continuation ratios \((0.11,0.08,0.06)\) were used for
the nonlinear solve unless a diagnostic run is explicitly described.  The
direct \texttt{mumps} linear solver was used.

The runner was executed with verbosity level \texttt{-v 2}.  At this level the
terminal logs include accepted Newton steps and line-search diagnostics.  Level
\texttt{-v 3} additionally enables PETSc KSP residual monitors when an
iterative linear solver is selected.

\section{Diagnostics}

The main closeness metric is the relative \(L^2\) design mismatch
\[
E_{\mathrm{rel}}
 =
\frac{\|\rho_{\mathrm{final}}-\rho_{\mathrm{design}}\|_{L^2(\Omega)}}
     {\|\rho_{\mathrm{design}}\|_{L^2(\Omega)}}.
\]
The study also records the absolute \(L^2\) mismatch, the signed mass
difference, and the active-set Jaccard overlap
\[
J_{\mathrm{active}}
=
\frac{|A_{\mathrm{final}}\cap A_{\mathrm{design}}|}
     {|A_{\mathrm{final}}\cup A_{\mathrm{design}}|},
\qquad
A = \{x : \rho(x) > 0.05\,\rho_{\mathrm{amp}}\}.
\]
The plateau Jaccard overlap was also written by the code, using the threshold
\(0.90\,\rho_{\mathrm{amp}}\), but it was zero for the completed runs and is
therefore not the discriminating metric in this sweep.

\section{Practical Lessons}

The clearest lesson is that the beta band should not be chosen as a copy of the
alpha band.  For the useful alpha window found on the coarse mesh,
\[
\alpha=(0.40,0.50),\qquad \alpha_c=0.45,\qquad \alpha_w=0.10,
\]
the successful beta bands are shifted substantially upward in the
\(\phi_{\mathrm{design}}\) scale.  The useful parameterization is
\[
\beta_w=\gamma\alpha_w,\qquad
\beta_c=\alpha_c+\delta,\qquad
\beta=(\beta_c-\beta_w/2,\ \beta_c+\beta_w/2).
\]

For relative \(L^2\) agreement with the torsion-designed density, the best
tested point was
\[
\gamma=1.50,\qquad \delta=0.425,\qquad \beta=(0.8000,0.9500),
\]
with \(E_{\mathrm{rel}}=0.384392\).  This point is close to the high-beta
robustness boundary: \(\gamma=1.50,\delta=0.45\) timed out.  A more robust and
geometrically balanced choice is
\[
\gamma=1.375,\qquad \delta=0.375,\qquad \beta=(0.7562,0.8938),
\]
which has \(E_{\mathrm{rel}}=0.409392\), \(J_{\mathrm{active}}=0.666506\), and
nearly zero mass mismatch.  If active-band overlap is the main objective, the
best tested band was
\[
\gamma=0.875,\qquad \delta=0.20,\qquad \beta=(0.6062,0.6937),
\]
with \(J_{\mathrm{active}}=0.757761\).

Thus the observed correlation is:
\[
\beta_c \approx \alpha_c + 0.35\text{--}0.425,\qquad
\beta_w \approx 1.25\text{--}1.50\,\alpha_w
\]
for relative-density matching, with smaller \(\gamma\) and smaller positive
\(\delta\) favored only when active-set overlap is prioritized.  Negative
\(\delta\) values were explicitly tested; they did not improve the match and
introduced several non-robust pockets for wider beta bands.

\subsection{Recommended Run Parameters}

For production torsion-initialized Newton runs in this study, use the classical
\(\phi_{\mathrm{design}}\)-scaled nonlinear window,
\[
c_{1,\Phi}=\beta_{\Phi,1}\max(\phi_{\mathrm{design}}),\qquad
c_{2,\Phi}=\beta_{\Phi,2}\max(\phi_{\mathrm{design}}),\qquad
\epsilon_{\Phi}=r_{\epsilon}(c_{2,\Phi}-c_{1,\Phi}).
\]
The best balanced dense-sweep choice is
\[
\alpha_{T,1}=0.40,\qquad
\alpha_{T,2}=0.50,\qquad
\beta_{\Phi,1}=0.75625,\qquad
\beta_{\Phi,2}=0.89375.
\]
In width/shift variables this same recommendation is
\[
\gamma=1.375,\qquad
\delta=0.375,\qquad
\beta_c=0.825,\qquad
\beta_w=0.1375.
\]
The more aggressive relative-\(L^2\) winner is
\[
\beta=(0.80,0.95),\qquad \gamma=1.50,\qquad \delta=0.425,
\]
but it is near the high-beta robustness boundary.  If active-set overlap is
the main objective, the best tested band was
\[
\beta=(0.60625,0.69375),\qquad \gamma=0.875,\qquad \delta=0.20.
\]

The runner also supports a torsion-scaled diagnostic window,
\[
w_T=(\alpha_{T,2}-\alpha_{T,1})\max(T),\qquad
c_{1,\Phi}=\alpha_{T,1}\max(T)+s\,w_T,\qquad
c_{2,\Phi}=c_{1,\Phi}+g\,w_T,
\]
through \texttt{--phi-window-source torsion},
\texttt{--phi-window-torsion-shift-scale}, and
\texttt{--phi-window-torsion-width-scale}.  This is not the recommended
production form for the present smooth-star setup: the tested \(p=5\)
torsion-scaled windows near \(\alpha=(0.40,0.50)\) placed \(c_{1,\Phi}\) far
above \(\max(\phi_{\mathrm{design}})\), leaving the nonlinear source
essentially empty and producing nonconverged zero-density final states.

\section{Initial Fine-Mesh Results}

\begin{table}[h]
\centering
\small
\begin{tabular}{@{}>{\raggedright\arraybackslash}p{0.31\linewidth}
                S[table-format=1.6]
                S[table-format=1.6]
                S[table-format=1.6]
                S[table-format=-1.6]
                S[table-format=2.2]@{}}
\toprule
Parameters & {$E_{\mathrm{rel}}$} & {$L^2$ diff.} & {$J_{\mathrm{active}}$} & {Mass diff.} & {Time (s)} \\
\midrule
\(\alpha=(0.45,0.55)\), \(\beta=(0.60,0.70)\) & 1.101402 & 0.834840 & 0.513855 & -0.031767 & 66.61 \\
\(\alpha=(0.50,0.60)\), \(\beta=(0.60,0.70)\) & 1.226505 & 0.916113 & 0.408180 & -0.008784 & 69.55 \\
\(\alpha=(0.55,0.65)\), \(\beta=(0.60,0.70)\) & 1.295088 & 0.956198 & 0.308372 & 0.009689 & 60.63 \\
\(\alpha=(0.60,0.70)\), \(\beta=(0.60,0.70)\) & 1.328580 & 0.972271 & 0.241929 & 0.024510 & 52.02 \\
\(\alpha=(0.65,0.75)\), \(\beta=(0.60,0.70)\) & 1.345975 & 0.978598 & 0.201530 & 0.035878 & 57.13 \\
\(\alpha=(0.60,0.70)\), \(\beta=(0.65,0.80)\) & 1.443650 & 1.056480 & 0.162580 & 0.226901 & 75.15 \\
\(\alpha=(0.60,0.70)\), \(\beta=(0.60,0.75)\) & 1.480901 & 1.083741 & 0.075095 & 0.278467 & 79.13 \\
\bottomrule
\end{tabular}
\caption{Completed fixed-mesh, \(p=4\), torsion-initialized Newton parameter runs.}
\end{table}

\section{Interrupted Runs}

Two low-\(\beta\) runs were stopped because their line-search histories showed
poor robustness and very low active-set overlap:

\begin{center}
\small
\begin{tabular}{@{}lrrrrrr@{}}
\toprule
Parameters & Stage & Residual & \(\alpha_{\mathrm{LS}}\) & BT & \(E_{\mathrm{rel}}\) & \(J_{\mathrm{active}}\) \\
\midrule
\(\alpha=(0.60,0.70)\), \(\beta=(0.50,0.65)\) & \(i_\epsilon=1,k=2\) & \(4.642694{\times}10^{-3}\) & \(2.384{\times}10^{-7}\) & 22 & 1.435762 & 0.004967 \\
\(\alpha=(0.60,0.70)\), \(\beta=(0.55,0.70)\) & \(i_\epsilon=0,k=18\) & \(3.521832{\times}10^{-3}\) & \(6.250{\times}10^{-2}\) & 4 & 1.430719 & 0.027175 \\
\bottomrule
\end{tabular}
\end{center}

\section{Initial Fine-Mesh Trend}

With \(\beta=(0.60,0.70)\) fixed, lowering the torsion window improved the
final agreement monotonically over the tested range.  The relative mismatch
dropped from \(1.345975\) at \(\alpha=(0.65,0.75)\) to \(1.101402\) at
\(\alpha=(0.45,0.55)\).  At the same time, the active-set Jaccard overlap rose
from \(0.201530\) to \(0.513855\).

With \(\alpha=(0.60,0.70)\) fixed, shifting the beta window upward worsened the
agreement.  The default \(\beta=(0.60,0.70)\) had
\(E_{\mathrm{rel}}=1.328580\), while \(\beta=(0.60,0.75)\) and
\(\beta=(0.65,0.80)\) increased the mismatch to \(1.480901\) and \(1.443650\),
respectively.  Shifting beta downward was not robust in this setup; the
line-search logs show small accepted steps, repeated backtracking, and active
overlap near zero.

\section{Initial Fine-Mesh Takeaway}

In the first fine-mesh bracket, before the wider coarse-mesh sweeps were run,
the closest design band to the final equilibrium band was obtained with
\[
\alpha_{T,1}=0.45,\qquad
\alpha_{T,2}=0.55,\qquad
\beta_{\Phi,1}=0.60,\qquad
\beta_{\Phi,2}=0.70.
\]
This was not the final recommendation.  It only established that the alpha
window needed to move below the original \((0.60,0.70)\) range.  The later
coarse-mesh sweeps refined this to \(\alpha=(0.40,0.50)\), then showed that the
beta band should be shifted upward relative to that alpha band.

\section{Artifacts}

The parsed completed-run table is available at
\[
\texttt{docs/strategyA\_band\_parameter\_study/results.csv}.
\]
The interrupted-case diagnostics are available at
\[
\texttt{docs/strategyA\_band\_parameter\_study/interrupted\_cases.csv}.
\]
Full verbose terminal logs are under
\[
\texttt{run\_outputs/strategyA\_band\_study\_20260712/runs/}.
\]

\section{Appended Coarse-Mesh Plotted Checks}

Additional checks were run on a coarser fixed mesh to make final-frame plotted
runs cheaper:
\[
\texttt{run\_outputs/strategyA\_band\_study\_20260712/fixed\_mesh/smooth\_star\_h010\_n220.msh}.
\]
This mesh has \(9{,}110\) triangles and \(73{,}321\) scalar degrees of freedom
for \(p=4\).  Each run enabled the final PyVista plot through
\texttt{--plot --plot-off-screen --save-frames --plot-final}.  Every completed
run saved exactly one final PNG frame.  The frame paths are listed in
\[
\texttt{docs/strategyA\_band\_parameter\_study/coarse\_final\_frame\_manifest.csv},
\]
and copied images are stored under
\[
\texttt{docs/strategyA\_band\_parameter\_study/final\_frames\_coarse/}.
\]
Because this mesh differs from the earlier \(16{,}238\)-triangle mesh, the
following values should be compared within this block rather than interpreted
as strict replacements for the fine-mesh values.

\begin{table}[h]
\centering
\scriptsize
\begin{tabular}{@{}>{\raggedright\arraybackslash}p{0.33\linewidth}
                r
                S[table-format=1.6]
                S[table-format=1.6]
                S[table-format=1.6]
                S[table-format=-1.6]
                r@{}}
\toprule
Parameters & {nt} & {$E_{\mathrm{rel}}$} & {$L^2$ diff.} & {$J_{\mathrm{active}}$} & {Mass diff.} & {Frames} \\
\midrule
\(\alpha=(0.40,0.50)\), \(\beta=(0.675,0.775)\) & 9110 & 0.484019 & 0.373529 & 0.582366 & -0.113546 & 1 \\
\(\alpha=(0.40,0.50)\), \(\beta=(0.65,0.75)\) & 9110 & 0.484042 & 0.373547 & 0.743608 & -0.096587 & 1 \\
\(\alpha=(0.375,0.475)\), \(\beta=(0.65,0.75)\) & 9110 & 0.497609 & 0.387993 & 0.631468 & -0.112520 & 1 \\
\(\alpha=(0.425,0.525)\), \(\beta=(0.65,0.75)\) & 9110 & 0.587736 & 0.449318 & 0.772225 & -0.082539 & 1 \\
\(\alpha=(0.35,0.45)\), \(\beta=(0.60,0.70)\) & 9110 & 0.694458 & 0.547477 & 0.659600 & -0.094118 & 1 \\
\(\alpha=(0.40,0.50)\), \(\beta=(0.625,0.725)\) & 9110 & 0.704700 & 0.543834 & 0.709753 & -0.078515 & 1 \\
\(\alpha=(0.45,0.55)\), \(\beta=(0.65,0.75)\) & 9110 & 0.730655 & 0.553824 & 0.747043 & -0.069589 & 1 \\
\(\alpha=(0.40,0.50)\), \(\beta=(0.60,0.70)\) & 9110 & 0.904823 & 0.698274 & 0.616926 & -0.059942 & 1 \\
\(\alpha=(0.35,0.45)\), \(\beta=(0.55,0.65)\) & 9110 & 0.959001 & 0.756030 & 0.520284 & -0.057169 & 1 \\
\(\alpha=(0.425,0.525)\), \(\beta=(0.60,0.70)\) & 9110 & 1.012675 & 0.774179 & 0.562224 & -0.045469 & 1 \\
\(\alpha=(0.45,0.55)\), \(\beta=(0.60,0.70)\) & 9110 & 1.101468 & 0.834893 & 0.509646 & -0.031623 & 1 \\
\(\alpha=(0.40,0.50)\), \(\beta=(0.55,0.65)\) & 9110 & 1.127332 & 0.869989 & 0.406804 & -0.020495 & 1 \\
\(\alpha=(0.45,0.55)\), \(\beta=(0.55,0.65)\) & 9110 & 1.251967 & 0.948969 & 0.299070 & 0.009687 & 1 \\
\bottomrule
\end{tabular}
\caption{Appended coarse-mesh checks with one saved final PyVista frame per run.}
\end{table}

The two best relative-\(L^2\) cases are nearly tied:
\[
\begin{aligned}
\alpha=(0.40,0.50),\ \beta=(0.675,0.775)
&:\quad E_{\mathrm{rel}}=0.484019,\quad J_{\mathrm{active}}=0.582366,\\
\alpha=(0.40,0.50),\ \beta=(0.65,0.75)
&:\quad E_{\mathrm{rel}}=0.484042,\quad J_{\mathrm{active}}=0.743608.
\end{aligned}
\]
The relative mismatch difference is only \(2.3\times 10^{-5}\), while the
active-band overlap is much better for \(\beta=(0.65,0.75)\).  On this
coarse plotted mesh, the more balanced choice is therefore
\[
\alpha_{T,1}=0.40,\qquad
\alpha_{T,2}=0.50,\qquad
\beta_{\Phi,1}=0.65,\qquad
\beta_{\Phi,2}=0.75.
\]
This appended family changes the earlier interpretation: once the torsion band
is moved down to the \((0.40,0.50)\) range, the preferred beta window shifts
upward rather than staying at \((0.60,0.70)\).  The next useful check is to
repeat the two leading coarse candidates on the finer \(16{,}238\)-triangle
mesh.

\section{Appended Extreme-Range Checks}

A further sweep tested windows closer to zero and closer to one on the same
coarse mesh,
\[
\texttt{run\_outputs/strategyA\_band\_study\_20260712/fixed\_mesh/smooth\_star\_h010\_n220.msh},
\]
with \(9{,}110\) triangles, \(73{,}321\) scalar degrees of freedom, and \(p=4\).
Each completed run saved one final PyVista PNG.  These copied final frames are
stored in
\[
\texttt{docs/strategyA\_band\_parameter\_study/final\_frames\_extremes/},
\]
with metadata in
\[
\texttt{docs/strategyA\_band\_parameter\_study/coarse\_extreme\_frame\_manifest.csv}.
\]

\begin{table}[h]
\centering
\scriptsize
\begin{tabular}{@{}>{\raggedright\arraybackslash}p{0.34\linewidth}
                r
                S[table-format=1.6]
                S[table-format=1.6]
                S[table-format=1.6]
                S[table-format=-1.6]
                r@{}}
\toprule
Parameters & {nt} & {$E_{\mathrm{rel}}$} & {$L^2$ diff.} & {$J_{\mathrm{active}}$} & {Mass diff.} & {Frames} \\
\midrule
\(\alpha=(0.70,0.80)\), \(\beta=(0.70,0.80)\) & 9110 & 0.826071 & 0.597888 & 0.816894 & -0.054066 & 1 \\
\(\alpha=(0.05,0.15)\), \(\beta=(0.05,0.15)\) & 9110 & 0.971700 & 0.904201 & 0.494615 & -0.386779 & 1 \\
\(\alpha=(0.10,0.20)\), \(\beta=(0.10,0.20)\) & 9110 & 1.105575 & 1.000852 & 0.324523 & -0.240512 & 1 \\
\(\alpha=(0.20,0.30)\), \(\beta=(0.20,0.30)\) & 9110 & 1.224105 & 1.045609 & 0.191160 & -0.068719 & 1 \\
\(\alpha=(0.80,0.90)\), \(\beta=(0.40,0.50)\) & 9110 & 1.520418 & 1.095090 & 0.000000 & 0.307132 & 1 \\
\bottomrule
\end{tabular}
\caption{Completed extreme-window checks.}
\end{table}

The following cases did not reach final convergence within the 210-second guard:
\[
\begin{array}{lll}
\alpha=(0.80,0.90), & \beta=(0.80,0.90), & \text{timeout},\\
\alpha=(0.90,0.98), & \beta=(0.90,0.98), & \text{timeout},\\
\alpha=(0.10,0.20), & \beta=(0.65,0.75), & \text{timeout},\\
\alpha=(0.20,0.30), & \beta=(0.65,0.75), & \text{timeout},\\
\alpha=(0.40,0.50), & \beta=(0.85,0.95), & \text{timeout}.
\end{array}
\]
For these timeout cases, a separate one-step diagnostic pass was run with
\(\epsilon\)-ratio \(0.11\) and \(\texttt{max-it}=1\).  The resulting PNGs are
nonconverged diagnostic states rather than final equilibria.  They are stored in
\[
\texttt{docs/strategyA\_band\_parameter\_study/diagnostic\_frames\_extreme\_timeouts/},
\]
with metadata in
\[
\texttt{docs/strategyA\_band\_parameter\_study/coarse\_extreme\_diagnostic\_frames.csv}.
\]

The extreme checks indicate that very low diagonal windows can converge, but
they are not competitive with the middle-window family and have large negative
mass differences.  The moderate high diagonal window
\(\alpha=(0.70,0.80)\), \(\beta=(0.70,0.80)\) converges and has strong active
overlap, but its relative mismatch remains worse than the best middle-window
candidate.  Near-one diagonal windows are slow or non-robust under the guard.
Low alpha with high beta, \(\alpha=(0.10,0.30)\), \(\beta=(0.65,0.75)\), is also
non-robust, and pushing beta to \((0.85,0.95)\) with the good alpha window times
out.  Thus the best practical coarse-mesh candidate remains
\[
\alpha=(0.40,0.50),\qquad \beta=(0.65,0.75).
\]

\section{Appended Beta Width/Shift Study}

The next sweep fixed the alpha window at
\[
\alpha_1=0.40,\qquad \alpha_2=0.50.
\]
The beta band was parameterized by a width multiplier \(\gamma\) and a center
shift \(\delta\):
\[
\alpha_w=\alpha_2-\alpha_1=0.10,\qquad
\alpha_c=\frac{\alpha_1+\alpha_2}{2}=0.45,
\]
\[
\beta_w=\gamma \alpha_w,\qquad
\beta_c=\alpha_c+\delta,\qquad
\beta_1=\beta_c-\frac{\beta_w}{2},\qquad
\beta_2=\beta_c+\frac{\beta_w}{2}.
\]
The first grid used
\[
\gamma\in\{0.50,0.75,1.00,1.25,1.50\},\qquad
\delta\in\{0.20,0.25,0.30\}.
\]
The best point was internal, so a local refinement used
\[
\gamma\in\{0.90,1.00,1.10\},\qquad
\delta\in\{0.235,0.250,0.265\}.
\]
All 23 runs completed on the same coarse mesh with \(9{,}110\) triangles and
\(73{,}321\) degrees of freedom, and every run saved one final PyVista PNG.
The copied figures are in
\[
\texttt{docs/strategyA\_band\_parameter\_study/final\_frames\_beta\_width\_shift/}.
\]

\begin{table}[h]
\centering
\scriptsize
\begin{tabular}{@{}cc>{\raggedright\arraybackslash}p{0.18\linewidth}
                S[table-format=1.6]
                S[table-format=1.6]
                S[table-format=1.6]
                S[table-format=-1.6]
                r@{}}
\toprule
\(\gamma\) & \(\delta\) & Beta band & {$E_{\mathrm{rel}}$} & {$L^2$ diff.} & {$J_{\mathrm{active}}$} & {Mass diff.} & {Frames} \\
\midrule
1.00 & 0.265 & \((0.665,0.765)\) & 0.443000 & 0.341873 & 0.648114 & -0.107128 & 1 \\
1.00 & 0.250 & \((0.650,0.750)\) & 0.484042 & 0.373547 & 0.743608 & -0.096587 & 1 \\
0.90 & 0.235 & \((0.640,0.730)\) & 0.484171 & 0.373647 & 0.618766 & -0.140009 & 1 \\
0.90 & 0.250 & \((0.655,0.745)\) & 0.579720 & 0.447384 & 0.513072 & -0.149851 & 1 \\
1.00 & 0.235 & \((0.635,0.735)\) & 0.611210 & 0.471686 & 0.735865 & -0.086059 & 1 \\
1.10 & 0.265 & \((0.660,0.770)\) & 0.627685 & 0.484399 & 0.725000 & -0.055101 & 1 \\
1.00 & 0.300 & \((0.700,0.800)\) & 0.669046 & 0.516319 & 0.431572 & -0.130029 & 1 \\
0.75 & 0.200 & \((0.6125,0.6875)\) & 0.676799 & 0.522302 & 0.455162 & -0.202082 & 1 \\
0.90 & 0.265 & \((0.670,0.760)\) & 0.697331 & 0.538147 & 0.420288 & -0.159394 & 1 \\
1.25 & 0.300 & \((0.6875,0.8125)\) & 0.712958 & 0.550207 & 0.692189 & -0.006361 & 1 \\
\bottomrule
\end{tabular}
\caption{Top beta width/shift results for fixed \(\alpha=(0.40,0.50)\).}
\end{table}

The best relative density mismatch is obtained at
\[
\gamma=1.00,\qquad \delta=0.265,\qquad \beta=(0.665,0.765),
\]
with \(E_{\mathrm{rel}}=0.443000\).  The previous
\(\beta=(0.650,0.750)\) remains the better active-overlap compromise:
\[
E_{\mathrm{rel}}=0.484042,\qquad J_{\mathrm{active}}=0.743608,
\]
compared with \(J_{\mathrm{active}}=0.648114\) for the relative-\(L^2\) best.
Very narrow beta widths, such as \(\gamma=0.50\), produced poor or zero active
overlap, while wider beta widths near \(\gamma=1.50\) increased mass and
relative density mismatch.

Thus, for relative-\(L^2\) closeness to the torsion-designed density, the best
tested beta band for fixed \(\alpha=(0.40,0.50)\) is
\[
\beta=(0.665,0.765).
\]
For a more geometric active-band match, \(\beta=(0.650,0.750)\) remains the
more balanced choice.

\section{Dense Beta Width/Shift Sweep With Negative Delta}

The previous beta width/shift sweep was expanded to include negative shifts and
many more samples.  The alpha window remained fixed at
\[
\alpha=(0.40,0.50).
\]
The dense grid used
\[
\gamma\in\{0.50,0.625,0.75,0.875,1.00,1.125,1.25,1.375,1.50\},
\]
and
\[
\delta\in\{-0.25,-0.20,\ldots,0.35\}.
\]
Because the best relative mismatch appeared on the upper shift boundary, the
high-shift side was extended with
\[
\gamma\in\{0.875,1.00,1.125,1.25,1.375,1.50\},\qquad
\delta\in\{0.375,0.40,0.425,0.45\}.
\]
This gives 141 attempted parameter pairs.  Of these, 122 completed with final
PyVista figures and 19 did not reach a final state within the guard.  Every
completed run has one copied final PNG in
\[
\texttt{docs/strategyA\_band\_parameter\_study/final\_frames\_beta\_width\_shift\_wide/}.
\]
The full CSV table and frame manifest are
\[
\texttt{coarse\_beta\_width\_shift\_wide\_results.csv}
\quad\text{and}\quad
\texttt{coarse\_beta\_width\_shift\_wide\_frame\_manifest.csv}.
\]
Heatmaps for relative mismatch, active Jaccard, mass difference, and timeouts
are stored as
\[
\texttt{coarse\_beta\_width\_shift\_wide\_rel\_heatmap.png},
\quad
\texttt{coarse\_beta\_width\_shift\_wide\_activej\_heatmap.png},
\]
\[
\texttt{coarse\_beta\_width\_shift\_wide\_massdiff\_heatmap.png},
\quad
\texttt{coarse\_beta\_width\_shift\_wide\_timeout\_map.png}.
\]

\begin{table}[h]
\centering
\scriptsize
\begin{tabular}{@{}cc>{\raggedright\arraybackslash}p{0.18\linewidth}
                S[table-format=1.6]
                S[table-format=1.6]
                S[table-format=-1.6]@{}}
\toprule
\(\gamma\) & \(\delta\) & Beta band & {$E_{\mathrm{rel}}$} & {$J_{\mathrm{active}}$} & {Mass diff.} \\
\midrule
1.500 & 0.425 & \((0.8000,0.9500)\) & 0.384392 & 0.583595 & 0.017897 \\
1.250 & 0.350 & \((0.7375,0.8625)\) & 0.408160 & 0.605462 & -0.041988 \\
1.375 & 0.375 & \((0.7562,0.8938)\) & 0.409392 & 0.666506 & -0.002054 \\
1.125 & 0.300 & \((0.6937,0.8063)\) & 0.425724 & 0.681023 & -0.067517 \\
1.375 & 0.400 & \((0.7812,0.9187)\) & 0.450885 & 0.536833 & -0.018853 \\
1.500 & 0.400 & \((0.7750,0.9250)\) & 0.469656 & 0.671479 & 0.035169 \\
1.000 & 0.250 & \((0.6500,0.7500)\) & 0.484042 & 0.743608 & -0.096587 \\
0.875 & 0.200 & \((0.6062,0.6937)\) & 0.548298 & 0.757761 & -0.129805 \\
\bottomrule
\end{tabular}
\caption{Top dense-grid cases by relative density mismatch.}
\end{table}

Negative \(\delta\) values were necessary to test, but they did not beat the
positive-shift families.  They generally produce worse relative density
mismatch and include several guarded timeouts for wider \(\gamma\).  For
relative \(L^2\) closeness to the torsion-designed density, the best tested
band is
\[
\gamma=1.50,\qquad \delta=0.425,\qquad \beta=(0.8000,0.9500).
\]
This point is near the high-shift robustness boundary:
\(\gamma=1.50,\delta=0.45\) timed out.  A more balanced choice is
\[
\gamma=1.375,\qquad \delta=0.375,\qquad \beta=(0.7562,0.8938),
\]
which has nearly the same relative mismatch as the second-best case, better
active overlap than the pure relative-\(L^2\) best, and almost zero mass
difference.  If active-band overlap is the primary criterion, the best tested
band is
\[
\gamma=0.875,\qquad \delta=0.20,\qquad \beta=(0.6062,0.6937).
\]

\end{document}
